\section{Conclusion}

This study set out to address a persistent gap in Kenya's livestock sector planning: the absence of a reproducible, data-driven framework linking herd demographics, production, and feed requirements across value chains. The resulting model resolves nine value chains, dairy cattle, beef cattle, camels, goats, sheep, poultry, pigs, rabbits, and apiculture, into age- and sex-structured cohorts advanced on a quarterly step from a 2019 baseline through 2041, translating mortality, culling, offtake, maturation, and reproduction into population, production, and feed demand trajectories for each chain.

The projections point to sustained, uneven growth across the sector. Ruminant numbers roughly triple over the horizon, led by beef cattle and goats, while non-ruminant growth is dominated by the divergence between a plateauing commercial poultry segment and a steadily expanding indigenous flock; production and feed demand follow correspondingly divergent paths by value chain. Beyond these headline trajectories, the model's cohort-level resolution surfaces the specific structural constraints, such as mortality concentrated in young cohorts, low off-take, and uneven reproductive performance, that drive them, offering a diagnostic view that aggregate, top-down projections cannot provide.

This combination of forecast and diagnosis is what gives the framework its practical value. Enterprise managers can use it to identify where herd performance is being lost and target culling, breeding, or feeding interventions accordingly; national and county planners can use it to move from static census snapshots to continuous, value-chain-specific investment and policy planning; and at the sector level, it offers a basis for anticipating feed and supply gaps before they become food security risks, particularly for the pastoralist and agro-pastoralist communities most exposed to them.

These contributions should be read alongside the limitations discussed above, chiefly the exclusion of environmental carrying capacity, gaps in Kenya-specific data for several value chains, and the omission of exogenous shocks such as disease and market fluctuations. Closing these gaps, through dedicated carrying-capacity work, targeted data collection, and extensions that incorporate economic and exogenous drivers, is the natural next step for turning this framework into a fuller decision-support tool for Kenya's livestock sector.
